% STATUS: written 2026-10-21, recalled results as is (one environment each, source kind and whole statement)
% STATUS: extended 2026-11-06 (Kawai, Guzzo and Ferreira 2025 read in full)
% Sources read in full: Fosner and Ilisevic 2008 (pp. 415--427), Ferreira, Guzzo and Ferreira 2020 (arXiv version, pp. 1--3), Majeed and Ozel 2014 (pp. 1881--1883); Kawai, Guzzo and Ferreira 2025 (pp. 1, 5--10, 12).
\section{Jordan triple, generalized Jordan and Jordan n-derivations}\label{sec:jordan-type:triple-and-generalized}

Section~\ref{sec:jordan-type:semiprime} showed that a Jordan derivation on a $2$-torsion free
semiprime ring is a derivation. Three families of maps change the rule $d(x^{2})=d(x)x+xd(x)$. Jordan triple derivations
state a rule for $xyx$ \cite[p.~416]{FosnerIlisevic2008}. Generalized Jordan derivations use a
second map \cite[p.~1]{FerreiraGuzzoFerreira2020}. The $(m,n)$-Jordan derivations of Vukman
weight the two terms \cite[p.~1881]{MajeedOzel2014}. This section takes each family in turn.

Fošner and Ilišević work with a ring $R$ and an $R$-bimodule $M$
\cite[p.~416]{FosnerIlisevic2008}. They define the following maps.

\begin{definition}[{\cite[Section~1, p.~416]{FosnerIlisevic2008}}]\label{def:jordan:triple-generalized}
Let $R$ be a ring and let $M$ be an $R$-bimodule. A Jordan triple derivation $\delta\colon R\to M$
is an additive mapping satisfying $\delta(xyx)=\delta(x)yx+x\delta(y)x+xy\delta(x)$ for all
$x,y\in R$. An additive mapping $E\colon R\to M$ is called a generalized derivation if there
exists a derivation $\delta\colon R\to M$ such that $E(xy)=E(x)y+x\delta(y)$ for all $x,y\in R$,
and it is called a generalized Jordan derivation if there is a Jordan derivation
$\delta\colon R\to M$ such that $E(x^{2})=E(x)x+x\delta(x)$ holds for all $x\in R$. An additive
mapping $E\colon R\to M$ is called a generalized Jordan triple derivation if there exists a
Jordan triple derivation $\delta\colon R\to M$ such that
$E(xyx)=E(x)yx+x\delta(y)x+xy\delta(x)$ for all $x,y\in R$.
\end{definition}

Every derivation is a Jordan derivation, and every generalized derivation is a generalized Jordan
derivation \cite[p.~416]{FosnerIlisevic2008}. Let the ring be $2$-torsion free. Then every Jordan
derivation is a Jordan triple derivation \cite[p.~416]{FosnerIlisevic2008}. Also every generalized
Jordan derivation is a generalized Jordan triple derivation with the same related map
\cite[p.~416]{FosnerIlisevic2008}. For semiprime rings
Fošner and Ilišević study a pair $(E,F)$ of additive maps \cite[p.~423]{FosnerIlisevic2008}.

\begin{theorem}[{\cite[Theorem~4.1, p.~424]{FosnerIlisevic2008}}]\label{thm:jordan:fosner-ilisevic-pair}
Let $R$ be a $2$-torsion free semiprime ring and let $E,F\colon R\to R$ be additive mappings.
Then the following conditions are mutually equivalent:
(a) for all $x,y\in R$, $E(xyx)=E(x)yx+xF(y)x+xyF(x)$ and $F(xyx)=F(x)yx+xF(y)x+xyF(x)$;
(b) for all $x,y,z\in R$, $E(xyz)=E(x)yz+xF(y)z+xyF(z)$ and $F(xyz)=F(x)yz+xF(y)z+xyF(z)$;
(c) for all $x,y\in R$, $E(xy)=E(x)y+xF(y)$ and $F(xy)=F(x)y+xF(y)$;
(d) for all $x\in R$, $E(x^{2})=E(x)x+xF(x)$ and $F(x^{2})=F(x)x+xF(x)$.
\end{theorem}

A direct consequence is the answer for generalized Jordan triple derivations
\cite[p.~424]{FosnerIlisevic2008}. It proves a conjecture posed in an earlier paper
\cite[p.~424]{FosnerIlisevic2008}.

\begin{corollary}[{\cite[Corollary~4.2, p.~424]{FosnerIlisevic2008}}]\label{cor:jordan:fosner-ilisevic-semiprime}
Let $R$ be a $2$-torsion free semiprime ring. Then any generalized Jordan triple derivation on $R$
is a generalized derivation.
\end{corollary}

Setting $F=0$ in Theorem~\ref{thm:jordan:fosner-ilisevic-pair} gives a statement about left
centralizers, which are the maps $T$ with $T(xy)=T(x)y$ \cite[p.~424]{FosnerIlisevic2008}.

\begin{corollary}[{\cite[Corollary~4.3, p.~424]{FosnerIlisevic2008}}]\label{cor:jordan:fosner-ilisevic-centralizer}
Let $R$ be a $2$-torsion free semiprime ring and let $T\colon R\to R$ be an additive mapping.
Then the following conditions are mutually equivalent:
(a) $T(xyx)=T(x)yx$ $(x,y\in R)$;
(b) $T(xyz)=T(x)yz$ $(x,y,z\in R)$;
(c) $T(xy)=T(x)y$ $(x,y\in R)$;
(d) $T(x^{2})=T(x)x$ $(x\in R)$.
\end{corollary}

Ferreira, Guzzo and Ferreira define a generalized Jordan derivation $\delta$ with a relating
Jordan derivation $\tau$ \cite[p.~1]{FerreiraGuzzoFerreira2020}. Jing and Lu introduced this
notion \cite[p.~2]{FerreiraGuzzoFerreira2020}. They stated the conjecture that every generalized Jordan
derivation on a $2$-torsion free semiprime ring is a generalized derivation
\cite[p.~2]{FerreiraGuzzoFerreira2020}. Ferreira, Guzzo and Ferreira prove it when the ring has a
nontrivial idempotent with mild conditions \cite[p.~2]{FerreiraGuzzoFerreira2020}. The conditions
follow the statement on the same page \cite[p.~2]{FerreiraGuzzoFerreira2020}.

\begin{theorem}[{\cite[Theorem~2.1, p.~2]{FerreiraGuzzoFerreira2020}}]\label{thm:jordan:ferreira-guzzo}
Let $R$ be a $2$-torsion free semiprime unity ring and $\delta$ be an additive generalized Jordan
derivation from $R$ into itself. If there exist an idempotent $e$ so that $e\neq0$, $e\neq1$ in $R$,
then $\delta$ is an additive generalized derivation.
\end{theorem}

The proof rests on the Peirce decomposition $R=R_{11}\oplus R_{12}\oplus R_{21}\oplus R_{22}$ with
respect to $e$ \cite[p.~2]{FerreiraGuzzoFerreira2020}. It also rests on a lemma of Jing and Lu
\cite[p.~3]{FerreiraGuzzoFerreira2020}.

\begin{lemma}[{\cite[Lemma~2.1, p.~3]{FerreiraGuzzoFerreira2020}}]\label{lem:jordan:jing-lu}
For all $a,b,c\in R$, the following statements hold:
(i) $\delta(ab+ba)=\delta(a)b+a\tau(b)+\delta(b)a+b\tau(a)$;
(ii) $\delta(aba)=\delta(a)ba+a\tau(b)a+ab\tau(a)$;
(iii) $\delta(abc+cba)=\delta(a)bc+a\tau(b)c+ab\tau(c)+\delta(c)ba+c\tau(b)a+cb\tau(a)$.
\end{lemma}

Majeed and Özel recall that Vukman introduced the third family \cite[p.~1881]{MajeedOzel2014}. Let $m\geq1$ and $n\geq1$ be fixed integers
with $m\neq n$, and let $A$ be an algebra \cite[p.~1881]{MajeedOzel2014}. A nonzero additive
mapping $\delta\colon A\to A$ with $(m+n)\delta(x^{2})=2m\delta(x)x+2nx\delta(x)$ for all $x\in A$
is called an $(m,n)$-Jordan derivation \cite[p.~1881]{MajeedOzel2014}.
Their first lemma shows how such a
map acts on idempotents.

\begin{lemma}[{\cite[Lemma~2.1, p.~1882]{MajeedOzel2014}}]\label{lem:jordan:majeed-ozel}
Let $A$ be a unital algebra. If $\delta$ is an $(m,n)$-Jordan derivation from $A$ into itself, then
for each idempotent $P\in A$, $(m+n)\delta(P)=2m\delta(P)P+2nP\delta(P)$.
\end{lemma}

Majeed and Özel then work in algebras of operators that leave the members of a commutative
subspace lattice invariant \cite[p.~1881]{MajeedOzel2014}.

The additive form of each of these notions has a common setting in axial
algebras \cite[p.~1]{KawaiGuzzoFerreira2025}. An axial algebra carries a
distinguished set of idempotents whose products follow a fusion rule
\cite[Definition~2.3, p.~5]{KawaiGuzzoFerreira2025}. A Jordan algebra with an
idempotent is an example, and so is the Norton--Sakuma algebra of type $2A$
\cite[Example~4.2, p.~8]{KawaiGuzzoFerreira2025}. Kawai, Guzzo and Ferreira
measure how far a multiplicative map is from being additive by a nullifying
function \cite[Definition~3.5, p.~6]{KawaiGuzzoFerreira2025}. This function is
symmetric in its finitely many arguments, and it vanishes on a repeated sum
exactly when it vanishes on the summands
\cite[Definition~3.5, p.~6]{KawaiGuzzoFerreira2025}. The authors prove that it
vanishes on a $J(\alpha)$-algebra under three conditions of Martindale type
\cite[Theorem~4.1, p.~8]{KawaiGuzzoFerreira2025}. Then four kinds of map are
additive. They are the $n$-multiplicative isomorphisms, the $n$-multiplicative
derivations, the surjective elementary maps, and, when the characteristic is not
$2$, the surjective Jordan elementary maps
\cite[Corollary~4.1, p.~10]{KawaiGuzzoFerreira2025}. The same conclusion holds
for $M(\alpha,\beta)$-axial algebras under five conditions
\cite[Theorem~5.1, p.~12]{KawaiGuzzoFerreira2025}.

On a $2$-torsion free semiprime ring the Jordan triple condition collapses to the
triple condition. Liu and Shiue show that every Jordan triple
$(\theta,\varphi)$-derivation is a $(\theta,\varphi)$-derivation, and deduce the
same for the generalized Jordan triple maps
\cite[Abstract]{LiuShiue2007JordanTriple}.

The comparison holds on a Lie ideal too. Hongan, Rehman and Al-Omary prove that on
a $2$-torsion free semiprime ring every Jordan triple derivation on a Lie ideal is
a derivation on it, and the same with "generalized" throughout
\cite[Abstract]{HonganRehmanAlOmary2011}.

Further work in the Jordan-type line, held in the collection:
\noindent\emph{Continuous linear maps on reflexive algebras behaving like Jordan left derivations at idempotent-product elements}~\cite{LibContinuousLinearMapsOn2013}.
\noindent\emph{Nonlinear Jordan derivable mappings of generalized matrix algebras by Lie product square-zero elements}~\cite{LibNonlinearJordanDerivableMappings2021}.
\noindent\emph{Jordan higher derivable maps on triangular algebras by commutative zero products}~\cite{LibJordanHigherDerivableMaps2016}.
